\documentclass[pdf, unicode, 10pt, notheorems, handout]{beamer}

%\mode<handout> {
%    \usepackage{pgfpages}
%    \setbeameroption{show notes}
%    \pgfpagesuselayout{2 on 1}[a4paper, border shrink=5mm]
%    \setbeamercolor{note page}{bg=white}
%    \setbeamercolor{note title}{bg=gray!10}
%    \setbeamercolor{note date}{fg=gray!10}
%}

\usepackage[utf8]{inputenc}
\usepackage[T2A]{fontenc}
\usepackage[russian, english]{babel}
\usepackage{amsmath, amssymb, amsfonts}
\usepackage{graphics}
\usepackage{mathtools}
\usepackage{tikz}
\usetikzlibrary{tikzmark, calc, fit}

\usefonttheme[onlymath]{serif}
\usetheme[numbers, totalnumbers, compress, nologo]{Statmod}
\setbeamerfont{institute}{size=\normalsize}
\setbeamertemplate{navigation symbols}{}

\graphicspath{{../img}}

\input{letters_series_mathbb}
\setbeamercolor{bluetext_color}{fg=blue}
\newcommand{\bluetext}[1]{{\usebeamercolor[fg]{bluetext_color}#1}}
\setbeamercolor{redtext_color}{fg=red}
\newcommand{\redtext}[1]{{\usebeamercolor[fg]{redtext_color}#1}}
\newcommand{\iu}{\mathbf{i}}
\newcommand{\Tp}{\mathrm{T}}
\newcommand{\Hp}{\mathrm{H}}
\newcommand{\rank}{\operatorname{rank}}

\newtheorem{theorem}{Теорема}
\newtheorem{statement}{Утверждение}
\newtheorem{definition}{Определение}
\newtheorem{remark}{Замечание}

\title{Тензорный анализ сингулярного спектра}

\author{Хромов Никита Андреевич, гр.24.M22-мм}

\institute[Санкт-Петербургский Государственный Университет]{
  Санкт-Петербургский государственный университет\\
  Математическое моделирование, программирование и \\искусственный интеллект\\
  Вычислительная стохастика и статистические модели\\
  \vspace{0.7cm}
Научный руководитель:  д.ф.-м.н. Голяндина Н.Э.}
%   \\Рецензент: к.ф.-м.н. Усевич К.Д.}

\date{Санкт-Петербург,\\ \number\year}

\begin{document}

\begin{frame}[plain, noframenumbering]
  \titlepage
\end{frame}

\section{Введение}\label{sec:intro}
\begin{frame}{Введение}
  $\tX^{(m)} = \left(x_1^{(m)}, x_2^{(m)}, \ldots,
  x_N^{(m)}\right)^{\rmT}$, $x_i \in \bbC$ ---
  временной ряд\\
  $\tX = \left(\tX^{(1)} : \tX^{(2)} : \ldots : \tX^{(M)}\right)$ ---
  $M$-канальный временной ряд\\ \vspace{0.2cm}
  $\tX = \tT + \tP + \tE$ \\
  $\tT$ --- медленно меняющаяся составляющая (тренд)\\
  $\tP$ --- периодическая составляющая (сезонность)\\
  $\tE$ --- случайная составляющая (шум)\\
  \vspace{0.2cm}

  \bluetext{Возможные задачи:}
  \begin{itemize}
    \item Выделение сигнала из ряда: нахождение $\tS = \tT + \tP$
    \item Разделение сигнала: нахождение составляющих $\tT$ и $\tP$
    \item Оценка параметров сигнала: определение периода и скорости
      изменения компонент сигнала
  \end{itemize}
  \smallskip

  Возможный метод решения: Singular Spectrum Analysis (\bluetext{SSA})

  %   [Broomhead, King (1986a)], [Golyandina, Nekrutkin, Zhigljavsky (2001)],\\
  %   и его многомерное расширение Multivariate SSA (\bluetext{MSSA})

  %   [Broomhead, King (1986b)]

  \smallskip

  \redtext{Цель:} сравнение и исследование свойств тензорных
  расширений методов SSA и MSSA с точки зрения точности выделения
  сигнала, разделения компонент и оценки параметров.
\end{frame}

\section{SSA}

\begin{frame}{Алгоиртм SSA}
  $\tX = \sum_{m=1}^{M}\tS_m + \tE$ --- временной ряд (один
  канал),\\ \vspace{0.1cm}
  $\tS \stackrel{\mathrm{def}}{=}\sum_{m=1}^{M}\tS_m$ ---
  сигнал.\\ \vspace{0.2cm}

  \bluetext{Параметры:} $L < N$ --- длина окна, $K
  \stackrel{\mathrm{def}}{=} N - L + 1$,\\
  $R$ --- число элементов разложения, относимых к сигналу,\\
  $\gI_1,\ldots, \gI_J \subseteq \{1, 2, \ldots, R\}$, $\gI_i \cap
  \gI_j = \varnothing$ --- наборы индексов,
  относимых к компонентам сигнала.
  \vspace{0.1cm}

  \bluetext{Алгоритм SSA для разделения компонент сигнала}
  \begin{enumerate}
    \item \bluetext{Вложение}
      $\tX \overset{L}{\longmapsto} \calT_L(\tX) = \bfX\in \bbC^{L\times K}$ ---
      траекторная матрица,

    \item \bluetext{Разложение} $\operatorname{SVD}(\bfX) =
      \sum_{i=1}^{d} \sqrt{\lambda}_{i}
      U_i V_i^{\Hp}, \quad d \leqslant \min(L,\, K)
      $
    \item \bluetext{Группировка} $\widetilde{\bfS}_j = \sum_{\color{red} i\in
      \gI_j}\sqrt{\lambda}_{i} U_i V_i^{\Hp}, \quad R \leqslant d
      $
    \item \bluetext{Восстановление} усреднение $\widetilde{\bfS}_j$
      вдоль антидиагоналей $l+k=\operatorname{const}$
  \end{enumerate}
  \vspace{0.1cm}

  \bluetext{Результат алгоритма} $\widetilde{\tS}_j$ --- оценки
  компонент $\tS_j$
\end{frame}

\begin{frame}{Траекторная матрица ряда}
  $\tX = (x_1,\, x_2,\, \dots,\, x_N)$, $L < N$, $K = N - L + 1$
  \vspace{0.4cm}\\
  \begin{definition}
    \[
      \bfX = \calT^{(L)}(\tX) =
      \begin{pmatrix}
        x_1                     & \tikzmarknode{A12}{x_2} &
        \tikzmarknode{A13}{x_3} & \ldots & x_K     \\
        \tikzmarknode{A21}{x_2} & x_3                     & x_4
        & \ldots & x_{K+1} \\
        \tikzmarknode{A31}{x_3} & x_4                     & x_5
        & \ldots & x_{K+2} \\
        \vdots                  & \vdots                  & \vdots
        & \ddots & \vdots  \\
        x_L                     & x_{L+1}                 & x_{L+2}
        & \ldots & x_N
      \end{pmatrix}
    \]
    \begin{tikzpicture}[remember picture,overlay]
      \draw[red] let \p1=($(A12)-(A21)$),\n1={atan2(\y1,\x1)} in
      node[rotate fit=\n1,fit=(A12) (A21),draw,rounded
      corners,inner sep=2pt]{};
      \draw[blue] let \p1=($(A13)-(A31)$),\n1={atan2(\y1,\x1)+2} in
      node[rotate fit=\n1,fit=(A13) (A31),draw,rounded
      corners,inner sep=2pt]{};
    \end{tikzpicture}
  \end{definition}
  \bigskip

  Шаг восстановления: $\widetilde{\tS} =
  \calT_L^{-1}\left(\Pi_{\calH}\left(\bfS_j\right)\right)$,\\
  $\Pi_{\calH}$ --- проектор на пространство ганкелевых матриц.
\end{frame}

\begin{frame}{SSA-Ранг Ряда}
  \begin{definition}
    Ряд $\tX$ имеет SSA-ранг $\rank(\tX)=r < N/2$, если ранг его $L$-траекторной
    матрицы равен $r$ для любой длины окна $L$ такой, что $r
    \leqslant \min(L, K)$.\\
    Если такое число $r$ существует, тогда ряд $\tX$ называется рядом
    конечного ранга.
  \end{definition}
  \vspace{0.5cm}
  Если сигнал $\tS$ является рядом конечного ранга, то в общем случае
  рекомендуется использовать $\rank(\tS)$ в качестве параметра $R$ в
  алгоритме SSA.
  \vspace{0.4cm}\\
  Примеры SSA-рангов
  \begin{itemize}
    \item ранг $\tS$ при $s_n = A \exp(\alpha n)$, $\alpha
      \in\bbC$, равен 1
    \item ранг $\tS$ при $s_n = A\sin(2\pi \omega n + \varphi)$,
      $0 < \omega < 1/2$, равен 2
  \end{itemize}
\end{frame}

\begin{frame}{Модель Сигнала}
  Ряд $\tS = (s_1,\, s_2,\, \ldots, s_N)$ считается сигналом, если
  траекторная матрица $\bfS = \calT^{(L)}(\tS)$ имеет неполный ранг
  ($\implies$ ряд имеет конечный ранг:
  $\operatorname{rank}(\tS)=r$)

  \bigskip

  Любой сигнал $\tS$ представим в форме конечной суммы:
  \[
    s_n = \sum_{j} p_j(n) \exp(\alpha_j n + \mathrm{i}2\pi
    \omega_j n),
  \]
  где $p_j(n)$ --- некоторый многочлен от $n$

  \bigskip

  Вещественный случай:
  \[
    s_n = \sum_{j} p_j(n) \exp(\alpha_j n)\sin(2\pi
    \omega_j n + \varphi_j),
  \]

  \bigskip
  ESPRIT --- близкий к SSA метод оценки степеней $\alpha_j$ и частот $\omega_j$

  %   [Roy et al. (1986)]
\end{frame}

\begin{frame}{Алгоритм ESPRIT}
  Входные данные: ряд $\tX$, длина окна $L$, число компонент сигнала $R$

  \bluetext{Алгоритм ESPRIT}
  \begin{enumerate}
    \item \bluetext{Вложение} $\bfX = \calT^{(L)}(\tX)$
    \item \bluetext{Разложение} $\displaystyle\bfX =
      \sum_{i=1}^{\operatorname{rank}\bfX}
      \sqrt{\lambda_i} U_i V_i^{\Hp}$,
      $\bfU_R = \left[U_1: U_2: \ldots: U_R\right]$
    \item \bluetext{Оценивание} Вычисление собственных чисел $z_j$ матрицы
      $\underline{\bfU}_R^{-}\overline{\bfU}_R$ \\ \smallskip
      Из равенства $z_j = \exp(\alpha_j + \mathrm{i}2\pi\omega_j)$ находятся
      параметры $\alpha_j$ и $\omega_j$.
  \end{enumerate}

  \bigskip

  Методы SSA, ESPRIT и их модификации называются методами,
  основанными на подпространстве сигнала.
\end{frame}

\section{Multi-Channel SSA}
\begin{frame}{MSSA: SSA для многоканальных рядов}
  %   [Broomhead, King (1986)]\\ \smallskip
  $\tX = \left(\tX^{(1)},\, \tX^{(2)}, \, \ldots,\,
  \tX^{(P)}\right)$ --- многоканальный ряд,
  $\tX^{(p)} = \left(x_1^{(p)}, \, x_2^{(p)},\,\ldots,\,
  x_N^{(p)}\right)$ --- каналы
  \vspace{0.4cm}

  Единственные отличия от алгоритма SSA --- шаг Вложения:
  \[
    \bfX = \calT_{\textrm{MSSA}}^{(L)}(\tX) = \left[\bfX^{(1)} :
    \bfX^{(2)}: \ldots : \bfX^{(P)}\right],\quad
    \bfX^{(p)} = \calT^{(L)}(\tX^{(p)}),
  \]
  и соответствующее изменение шага Восстановления.
  \vspace{0.4cm}

  Когда MSSA оказывается точнее, чем применение SSA к каждому каналу:
  \begin{itemize}
    \item Все каналы имеют ``похожую'' структуру
    \item Имеются ``поддерживающие'' каналы с меньшим уровнем шума
  \end{itemize}
\end{frame}

\section{Tensor SSA}
\begin{frame}{Переход к тензорам}
  \begin{table}
    \hspace*{-0.9cm}
    \begin{tabular}{rlllll}
      \bluetext{Basic SSA:} & Ряд $\tX$ & $\mapsto$ & Матрица
      $\mathbf{X}$ & $\mapsto$ & $\operatorname{SVD}(\mathbf{X})$ \\
      \bluetext{Tensor SSA:} & Ряд $\tX$ & $\mapsto$ & Тензор
      $\calX $ & $\mapsto$ & \hspace*{-0.18cm}
      $\underbrace{\operatorname{TD}}_{\mathclap{\text{~~~Тензорное
      Разложение}}}(\calX)$
    \end{tabular}
  \end{table}
  \vspace{0.4cm}

  Тензорные разложения, обобщающие SVD:
  \begin{itemize}
    \item \bluetext{Higher-Order SVD (HOSVD)}
    \item Canonical Polyadic Decomposition (CPD)
    \item T-SVD
    \item $(L_r, L_r, 1)$-Decomposition
  \end{itemize}
\end{frame}

\begin{frame}{Вложение одноканального ряда в Тензор}
  %   Considered in [Papy et al. (2005)]\\ \smallskip
  $\tX = (x_1,\, x_2,\, \ldots,\, x_N)$, $L, K < N - 1$, $M = N - L -
  K + 2 > 1$\\ \smallskip
  $\tX \mapsto \calT_{\mathrm{T-SSA}}^{(L, K)}(\tX) = \calX \in
  \bbC^{L\times K \times M}$:

  \begin{figure}[!ht]
    \centering
    \includegraphics[width=0.75\textwidth]{tens-injection-wide.pdf}
  \end{figure}
\end{frame}

\begin{frame}{Вложение многоканального ряда в Тензор}
  %   Considered in [Papy et al. (2005)]\\ \smallskip
  $\tX = \left(\tX^{(1)},\, \tX^{(2)},\, \ldots,\,
  \tX^{(M)}\right)$, $\tX^{(m)} = (x_1^{(m)},\, x_2^{(m)},\,
  \ldots,\, x_N^{(m)})$, $L < N$\\ \smallskip
  $\tX \mapsto \calT_{\mathrm{T-MSSA}}^{(L)}(\tX)
  = \calX \in \bbC^{L \times K \times M}$, $K = N - L + 1$

  \bigskip

  \begin{figure}[!ht]
    \centering
    \includegraphics[width=\textwidth]{mssa_injection_new.pdf}
  \end{figure}
\end{frame}

\begin{frame}{Некоторые тензорные разложения}
  В отличие от матричного случая, существует несколько различных
  определений тензорных рангов и соответствующих им обобщений SVD.

  \begin{definition}[Ранг тензора]
    Тензор $\calA$ имеет ранг 1, если существуют векторы $B$, $C$ и
    $D$ такие, что
    $\calA=B \circ C \circ D$, где $\circ$ обозначает тензорное
    произведение.\\ \medskip

    Тензор $\calA$ имеет ранг $R$, если его можно представить в виде суммы
    $R$ одноранговых тензоров: $\calA = \sum_{i=1}^{R} \calB_i$,
    $\rank(\calB_i)=1$, и такое $R$ минимально.
  \end{definition}

  \medskip

  \bluetext{Canonical Polyadic Decomposition (CPD)} --- представление
  тензора в виде суммы $R$ одноранговых тензоров.

  \smallskip

  \begin{itemize}
      % \item Considered for signal extraction in [Kouchaki, Sanei (2013)]
    \item Компоненты разложения не ортогональны
    \item Необходимо знать параметр $R$ до разложения
    \item Нет связи между SSA-рангом и рангом траекторного тензора
  \end{itemize}
\end{frame}

\begin{frame}{Некоторые тензорные разложения}
  \begin{definition}
    \bluetext{$n$-векторами} тензора $\calA$ называются векторы,
    получаемые из $\calA$ варьированием индекса $n$-го направления
    при фиксированных остальных (аналог строк и столбцов матрицы).\\ \medskip

    \bluetext{$n$-рангом} тензора $\calA$, обозначаемым $R_n =
    \rank_n(\calA)$, называется размерность линейного пространства,
    порождённого $n$-векторами.
  \end{definition}

  \medskip

  \bluetext{HOSVD}: $\calA = \sum_{l=1}^{R_1}\,
  \sum_{k=1}^{R_2}\, \sum_{m=1}^{R_3}\, z_{lkm} U_l^{(1)} \circ
  U_k^{(2)} \circ U_m^{(3)}$
  \begin{itemize}
      % \item Considered for signal parameter estimation in [Papy et
      %   al. (2005)], [Papy et al. (2009)]
    \item Компоненты разложения ортогональны
    \item Не требует заранее знать $n$-ранги тензора
    \item Была доказана связь между SSA-рангом сигналов и $n$-рангами
      их траекторных тензоров
  \end{itemize}
\end{frame}
\end{document}